\begin{tomtat}
\begin{itemize}
  \item Một \tn{attribution nối tầng}{serial attribution} lấy chính phương pháp attribution áp lên đầu
  ra của nó, như phương trình~\ref{eq:ch11-noi-tang} viết, và đó là nghĩa đen
  của nhan đề chương. Cấu trúc ấy có một khiếm khuyết mà bài báo gọi là rò rỉ
  tương tác và chứng minh trên một hàm hai biến có đáp án đúng biết trước:
  bảng~\ref{tab:ch11-tach-hieu-ung} cho thấy serial Shapley value gán
  $x_1 + \tfrac{1}{4}T$ cho một \tn{hiệu ứng riêng}{pure individual effect} lẽ ra bằng $x_1$, còn integrated
  Hessians gán $\tfrac{4}{9}T$ cho một hiệu ứng riêng lẽ ra bằng không. Đáp án
  đúng ở đó có được vì hàm và baseline được chọn sao cho nó có, cùng cách dựng
  ground truth mà chương~\ref{ch:chi-so-khong-do-do-trung-thuc} đọc, nên nó
  không nói được gì về một transformer.

  \item Định nghĩa~\ref{def:ch11-meta-attribution} sửa chỗ rò bằng cách đổi hàm
  giá trị của trò chơi hợp tác: người chơi vẫn là các đặc trưng, nhưng cái được
  chia không còn là đầu ra của mô hình mà là chính attribution bậc nhất. Giữ cố
  định đặc trưng $i$ rồi lấy giá trị Shapley của $j$ trên $d-1$ đặc trưng còn
  lại là chỗ khiến con số thu được có chiều, thứ mà các chỉ số theo tập không
  có. Định lý~\ref{thm:ch11-phan-tang-san-co} cho thấy các chỉ số tương tác sẵn
  có đã phân tách attribution theo đúng kiểu ấy từ trước, ở chỗ không ai phát
  biểu điều đó.

  \item \tn{Tính hiệu quả phân tầng}{hierarchical efficiency} của
  định nghĩa~\ref{def:ch11-hieu-qua-phan-tang}
  là tính đầy đủ của mệnh đề~\ref{prop:ch05-tinh-day-du} nâng lên một tầng: tổng
  các phần bằng cái được chia, chỉ khác cái được chia giờ là attribution chứ
  không phải đầu ra mô hình. Tiên đề ấy ràng buộc quan hệ giữa hai tầng và
  không nói gì về việc tầng dưới có đúng hay không. Meta-G$\times$I là ví dụ:
  nó hiệu quả phân tầng trong khi gradient nhân đầu vào nằm dưới nó cho tổng
  $x_1 + 3T$ trên một hàm có hiệu số đầu ra bằng $x_1 + T$.

  \item Ba thí nghiệm của bài lấy ba cái nền, và cả ba đều nằm ở phía tính hợp
  lý của cặp khái niệm mà mục~\ref{sec:ch01-do-trung-thuc-va-tinh-hop-ly} tách
  ra: cặp token do chính các tác giả gán nhãn, nhãn lớp của ImageNet-1k, và mặt
  nạ phân đoạn. Không cái nào trả lời câu hỏi mô hình đã dùng cái gì.
  Hình~\ref{fig:ch11-ba-tang} xếp ba tầng lại và đánh dấu bằng nét chấm hai
  tầng chưa có dụng cụ đo đã được kiểm định.

  \item Phần hạn chế của bài viết rằng một phép so sánh phương pháp cho kết
  luận chặt còn cần thêm phép đo độ trung thực, cụ thể là các đường cong
  deletion và insertion trên từng cặp token, và bài không thực hiện phép đo ấy.
  Hai đường cong đó thuộc các họ chỉ số mà mục~\ref{sec:ch08-ho-xoa-dan} đã đọc
  và chưa cái nào được kiểm định. Vì vậy vòng nghiên cứu ở đây chạy đủ một lượt
  mà không có bước loại bỏ: có chỗ yếu, có xây dựng mới sửa được chỗ yếu ấy, có
  phép so theo chỉ số do người đề xuất chọn, có công bố, và không có bước nào
  nói được rằng một phương pháp cũ nên thôi dùng.
\end{itemize}
\end{tomtat}

\cauhoi

\begin{cauhoids}
  \item Áp phương trình~\ref{eq:ch11-noi-tang} thêm một lần nữa, tức lấy
  attribution của attribution của attribution. Viết ra biểu thức thu được, cho
  biết nó gán một con số cho đối tượng nào, và nói xem
  định nghĩa~\ref{def:ch11-hieu-qua-phan-tang} có phát biểu được cho tầng thứ
  ba ấy không.

  \item Với hàm $f(x) = x_1 + x_1 x_2^2$ và baseline tại gốc, hãy đổi vai hai
  đặc trưng bằng cách xét $f(x) = x_2 + x_1^2 x_2$ và tính lại hai hiệu ứng
  riêng theo định nghĩa của bài. Dòng \emph{đáp án đúng} của
  bảng~\ref{tab:ch11-tach-hieu-ung} đổi thế nào, và điều đó cho biết gì về việc
  đáp án ấy đến từ đâu?

  \item Định nghĩa~\ref{def:ch11-hieu-qua-phan-tang} và
  mệnh đề~\ref{prop:ch05-tinh-day-du} có cùng dạng. Hãy nêu rõ ở mỗi chỗ, cái
  được chia là gì và các phần là gì. Sau đó cho biết một phương pháp thỏa cả
  hai thì đã được bảo đảm điều gì và chưa được bảo đảm điều gì.

  \item Ba thí nghiệm của mục~\ref{sec:ch11-ba-phep-do-va-ba-cai-nen} dùng ba
  cái nền khác nhau. Giả sử phải thiết kế một phép đo thứ tư cho
  meta-attribution mà nền của nó nằm ở phía độ trung thực. Phép đo ấy phải trả
  lời câu hỏi nào, và theo chương~\ref{ch:do-do-trung-thuc} thì dụng cụ nào
  hiện có để trả lời, kèm điều kiện nào chưa được xác lập cho dụng cụ ấy?

  \item Đọc phần hạn chế ở cuối phần chính của bài
  báo~\autocite{p23metagame}, gồm bốn phát biểu. Xếp chúng thành hai nhóm, nhóm
  nói về giới hạn của xây dựng lý thuyết và nhóm nói về giới hạn của bằng chứng
  thực nghiệm, rồi cho biết nhóm nào chương~\ref{ch:khoang-trong-nghien-cuu} sẽ
  cần tới khi tổng hợp khoảng trống nghiên cứu.
\end{cauhoids}
